\frfilename{mt115.tex}
\versiondate{21.7.05}
\copyrightdate{1994}

\def\chaptername{Ruang ukur}
\def\sectionname{Ukuran Lebesgue pada $\BbbR^r$}


\newsection{115}
\def\headlinesectionname{Ukuran Lebesgue pada $\eightBbb R^r$}

Sesudah gagasan-gagasan yang sangat abstrak dalam \S\S111-113, terdapat
kebutuhan mendesak akan contoh-contoh ruang ukur yang taktrivial. Contoh
yang jauh lebih penting daripada contoh-contoh lain adalah ruang-ruang
Euklides $\BbbR^r$ dengan ukuran Lebesgue, dan sekarang saya
akan menguraikan definisi ukuran-ukuran ini (115A-115E), beserta beberapa
sifat dasarnya. Kecuali pada satu titik (dalam bukti lema fundamental
115B), bagian ini tidak bergantung secara hakiki pada \S114; tetapi
bagaimanapun, kebanyakan mahasiswa yang menjumpai ukuran Lebesgue untuk
pertama kalinya akan merasa lebih mudah mempelajari dengan saksama kasus
berdimensi satu sebelum beranjak ke kasus berdimensi banyak.

\leader{115A}{Definisi (a)} Hampir di seluruh bagian ini
(pengecualiannya adalah bukti Lema 115B), $r$ akan menyatakan suatu
bilangan bulat tetap yang lebih besar daripada atau sama dengan $1$.
Saya akan menggunakan huruf Latin $a$, $b$, $c$, $d$, $x$, $y$ untuk
menyatakan anggota $\BbbR^r$, dan huruf Yunani untuk koordinatnya,
sehingga $a=(\alpha_1,\ldots,\alpha_r)$,
$b=(\beta_1,\ldots,\beta_r)$, $x=(\xi_1,\ldots,\xi_r)$.

\header{115Ab}{\bf (b)} Untuk keperluan bagian ini, sebuah {\bf
interval setengah terbuka} dalam $\BbbR^r$ adalah himpunan berbentuk
$\coint{a,b}=\{x:\alpha_i\le \xi_i<\beta_i\Forall i\le
r\}$, dengan $a$, $b\in\BbbR^r$.
Perhatikan bahwa saya mengizinkan $\beta_i\le \alpha_i$ dalam rumus ini;
jika hal ini terjadi untuk suatu $i$, maka $\coint{a,b}=\emptyset$.



\header{115Ac}{\bf (c)} Jika $I=\coint{a,b}\subseteq \BbbR^r$ adalah
interval setengah terbuka, maka
$I=\emptyset$ atau

\Centerline{$\alpha_i=\inf\{\xi_i:x\in I\}$,
\quad$\beta_i=\sup\{\xi_i:x\in I\}$}

\noindent untuk setiap $i\le r$; dalam kasus kedua, penyajian $I$
sebagai interval setengah terbuka bersifat unik.
Karena itu, kita dapat mendefinisikan {\bf volume berdimensi $r$}
$\lambda I$ dari suatu interval setengah terbuka $I$ dengan menetapkan

\Centerline{$\lambda\emptyset = 0$,
\quad $\lambda\coint{a,b}=\prod_{i=1}^r\beta_i-\alpha_i$ jika
$\alpha_i<\beta_i$ untuk setiap $i$.}



\leader{115B}{Lema} Jika $I\subseteq\BbbR^r$ adalah interval setengah
terbuka dan $\langle I_j\rangle_{j\in\Bbb N}$ adalah barisan interval
setengah terbuka yang menutupi $I$, maka
$\lambda I\le\sum_{j=0}^{\infty}\lambda I_j$.

\proof{ Bukti dilakukan dengan induksi pada $r$. Karena itu, hanya untuk
bukti ini, saya menuliskan $\lambda_r$ bagi fungsi yang didefinisikan
pada interval-interval setengah terbuka dalam $\BbbR^r$ melalui rumus
115Ac.

\medskip

{\bf (a)} Argumen untuk $r=1$, yang memulai induksi, serupa dengan
langkah induksi; tetapi alih-alih menetapkan kesepakatan yang sesuai
untuk membangun kasus trivial $r=0$, atau meminta Anda mengerjakan sendiri
rinciannya, saya merujuk Anda ke 114B, yang tepat merupakan kasus $r=1$.

\medskip

{\bf (b)} Untuk langkah induksi menuju $r+1$, dengan $r\ge 1$, ambil
interval setengah terbuka $I\subseteq\BbbR^{r+1}$ dan
$\sequence{j}{I_j}$ suatu barisan interval setengah terbuka yang
menutupi $I$. Jika $I=\emptyset$, tentu saja $\lambda_{r+1}
I=0\le\sum_{j=0}^{\infty}\lambda_{r+1}
I_j$. Jika tidak, nyatakan
$I$ sebagai $\coint{a,b}$, dengan $\alpha_i<\beta_i$ untuk $i\le r+1$,
dan setiap $I_j$ sebagai $\coint{a^{(j)},b^{(j)}}$. Tuliskan
$\zeta=\prod_{i=1}^r\beta_i-\alpha_i$, sehingga
$\lambda_{r+1}I=\zeta(\beta_{r+1}-\alpha_{r+1})$. Tetapkan
$\epsilon>0$. Untuk setiap $\xi\in \Bbb R$, misalkan $H_{\xi}$ adalah
setengah-ruang $\{x:\xi_{r+1}<\xi\}$, dan tinjau himpunan

\Centerline{$A=\{\xi:\alpha_{r+1}\le \xi\le \beta_{r+1},\,
\zeta(\xi-\alpha_{r+1})
\le(1+\epsilon)\sum_{j=0}^{\infty}\lambda_{r+1}(I_j\cap H_{\xi})\}$.}

\noindent (Perhatikan bahwa
$I_j\cap H_{\xi}
=\hbox{$\bigl[$}a^{(j)},\tilde b^{(j)}\hbox{$\bigr[$}$,
dengan $\tilde\beta^{(j)}_i=\beta^{(j)}_i$ untuk $i\le r$ dan
$\tilde\beta^{(j)}_{r+1}=\min(\beta^{(j)}_{r+1},\xi)$, sehingga
$\lambda_{r+1}(I_j\cap H_{\xi})$ selalu terdefinisi.)
Kita mempunyai $\alpha_{r+1}\in A$, karena

\Centerline{$\zeta(\alpha_{r+1}-\alpha_{r+1})=0
\le(1+\epsilon)\sum_{j=0}^{\infty}\lambda_{r+1}(I_j\cap H_{\alpha_{r+1}})$,}

\noindent dan tentu saja $A\subseteq[\alpha_{r+1},\beta_{r+1}]$,
sehingga $\gamma=\sup A$ terdefinisi dan
termasuk dalam $[\alpha_{r+1},\beta_{r+1}]$.

\medskip

{\bf (c)} Sekarang kita dapati bahwa $\gamma\in A$.

$$\eqalign{\Prf\ \zeta(\gamma-\alpha_{r+1})
&=\sup_{\xi\in A}\zeta(\xi-\alpha_{r+1})\cr
&\le(1+\epsilon)\sup_{\xi\in A}
  \sum_{j=0}^{\infty}\lambda_{r+1}(I_j\cap H_{\xi})
\le(1+\epsilon)\sum_{j=0}^{\infty}\lambda_{r+1}(I_j\cap H_{\gamma}).
\text{ \Qed}\cr}$$

\wheader{115B}{6}{2}{2}{12pt}

{\bf (d)} \Quer\ Andaikan, jika mungkin, bahwa
$\gamma<\beta_{r+1}$. Maka
$\gamma\in\coint{\alpha_{r+1},\beta_{r+1}}$. Tetapkan

\Centerline{$J=\{x:x\in\BbbR^r,\,(x,\gamma)\in I\}
=\coint{a',b'}$,}

\noindent dengan $a'=(\alpha_1,\ldots,\alpha_r)$,
$b'=(\beta_1,\ldots,\beta_r)$, dan untuk setiap $j\in\Bbb N$
tetapkan

\Centerline{$J_j=\{x:x\in\BbbR^r,\,(x,\gamma)\in I_j\}$.}

\noindent Karena $I\subseteq\bigcup_{j\in\Bbb N}I_j$, kita harus
mempunyai $J\subseteq\bigcup_{j\in\Bbb N}J_j$. Tentu saja, baik $J$
maupun setiap $J_j$ merupakan interval setengah terbuka dalam
$\BbbR^r$. (Inilah salah satu tempat yang memperlihatkan manfaat
menganggap himpunan kosong sebagai interval setengah terbuka.)
Menurut hipotesis induksi,
$\zeta=\lambda_rJ\le\sum_{j=0}^{\infty}\lambda_rJ_j$. Karena $\zeta>0$,
terdapat $m\in\Bbb N$ sedemikian sehingga
$\zeta\le(1+\epsilon)\sum_{j=0}^m\lambda_rJ_j$.
Sekarang, untuk setiap $j\le m$, berlaku $J_j=\emptyset$ atau
$\alpha^{(j)}_{r+1}\le\gamma<\beta^{(j)}_{r+1}$; tetapkan

\Centerline{$\xi=\min(\{\beta_{r+1}\}\cup\{\beta^{(j)}_{r+1}:j\le
m,\,J_j\ne\emptyset\})>\gamma$.}

\noindent Maka

\Centerline{$\lambda_{r+1}(I_j\cap H_{\xi})
\ge\lambda_{r+1}(I_j\cap H_{\gamma})+(\xi-\gamma)\lambda_rJ_j$}

\noindent untuk setiap $j\le m$ yang untuknya $J_j$ tak kosong, dan
karena itu untuk setiap $j$. Akibatnya,

$$\eqalign{\zeta(\xi-\alpha_{r+1})
&=\zeta(\gamma-\alpha_{r+1})+\zeta(\xi-\gamma)\cr
&\le(1+\epsilon)\sum_{j=0}^{\infty}\lambda_{r+1}(I_j\cap H_{\gamma})
+(1+\epsilon)(\xi-\gamma)\sum_{j=0}^m\lambda_rJ_j\cr
&\le(1+\epsilon)\sum_{j=m+1}^{\infty}\lambda_{r+1}(I_j\cap H_{\gamma})
+(1+\epsilon)\sum_{j=0}^m\lambda_{r+1}(I_j\cap H_{\xi})\cr
&\le(1+\epsilon)\sum_{j=0}^{\infty}\lambda_{r+1}(I_j\cap H_{\xi}),\cr}$$

\noindent dan $\xi\in A$, yang mustahil.\ \Bang

\medskip

{\bf (e)} Kita menyimpulkan bahwa $\gamma=\beta_{r+1}$, sehingga
$\beta_{r+1}\in A$ dan

\Centerline{$\lambda_{r+1}I=\zeta(\beta_{r+1}-\alpha_{r+1})
\le(1+\epsilon)\sum_{j=0}^{\infty}\lambda_{r+1}(I_j\cap H_{\beta_{r+1}})
\le(1+\epsilon)\sum_{j=0}^{\infty}\lambda_{r+1}I_j$.}

\noindent Karena $\epsilon$ sebarang,

\Centerline{$\lambda_{r+1}I\le\sum_{j=0}^{\infty}\lambda_{r+1}I_j$,}

\noindent sebagaimana diklaim.
}%end of proof of 115B

\cmmnt{\medskip
\noindent{\bf Catatan} Bukti ini cukup berat, dan tidak semua orang
menyusun uraian
\ifUSEnglish{serinci}\else\fi
ini untuk membuktikannya. Pendekatan yang barangkali lebih konvensional diuraikan
secara garis besar dalam 115Ya, dengan menggunakan teorema Heine--Borel
untuk mereduksi masalah menjadi persoalan tentang liputan berhingga,
kemudian (sangat sering) mengatakan bahwa persoalan tersebut trivial.
Saya tidak menggunakan metode ini, antara lain karena kita tidak
memerlukan teorema Heine--Borel di tempat lain dalam jilid ini (meskipun
kita pasti memerlukannya dalam Jilid 2, dan saya menuliskan buktinya
dalam 2A2F), dan antara lain karena saya tidak setuju bahwa lema tersebut
trivial ketika kita mempunyai barisan berhingga $I_0,\ldots,I_m$ yang
menutupi $I$. Saya mengundang Anda untuk mempertimbangkannya sendiri.
Menurut saya, setiap argumen yang saksama harus melibatkan induksi pada
dimensi, dan itulah yang saya berikan di sini. Tentu saja, penanganan
barisan tak hingga sepanjang pembuktian membuat pelacakan atas apa yang
sedang kita kerjakan sedikit lebih sulit, dan saya mencatat bahwa
sesungguhnya terdapat suatu langkah krusial yang mengharuskan pemotongan
barisan; yang saya maksud ialah rumus

\Centerline{$\xi=\min(\{\beta_{r+1}\}\cup\{\beta^{(j)}_{r+1}:
                                          j\le m,\,J_j\ne\emptyset\})$}

\noindent pada bagian (d) dari bukti. Kita jelas tidak dapat mengambil
$\xi=\inf\{\beta^{(j)}_{r+1}:j\in\Bbb N,\,J_j\ne\emptyset\}$, karena
nilai ini sangat mungkin sama dengan $\gamma$. Dengan demikian, saya
memerlukan suatu alasan untuk memotong barisan, dan alasan itu terdapat
dalam kalimat

\Centerline{Karena $\zeta>0$, terdapat $m\in\Bbb N$ sedemikian sehingga
$\zeta\le(1+\epsilon)\sum_{j=0}^m\lambda_rJ_j$.}

\noindent Langkah itulah alasan untuk memasukkan kelonggaran
$\epsilon$ ke dalam definisi himpunan $A$ pada awal bukti. Selain
modifikasi ini, struktur argumen tersebut dimaksudkan mencerminkan
struktur 114B; jadi saya berharap Anda dapat menggunakan rumus-rumus
yang lebih sederhana dalam 114B sebagai panduan di sini.
}%end of comment

\leader{115C}{Definisi} Sekarang, dan untuk seluruh sisa bagian ini,
definisikan $\theta:\Cal P(\BbbR^r)\to[0,\infty]$ dengan
menetapkan

$$\eqalign{\theta A=\inf\{\sum_{j=0}^{\infty}\lambda I_j:\langle
I_j\rangle_{j\in\Bbb N}
\text{ adalah barisan }&\text{interval setengah terbuka}\cr
&\text{sedemikian sehingga }
A\subseteq\bigcup_{j\in\Bbb N}I_j\}.\cr}$$


\cmmnt{\noindent Perhatikan bahwa setiap $A$ dapat ditutupi oleh suatu
barisan interval setengah terbuka --- misalnya,
$A\subseteq\bigcup_{n\in\Bbb N}\coint{-\tbf{n},\tbf{n}}$, dengan
$\tbf{n}=(n,n,\ldots,n)\in\BbbR^r$; sehingga jika kita menafsirkan
jumlah-jumlahnya dalam $[0,\infty]$, seperti dalam 112Bc di atas, kita
selalu mempunyai himpunan tak kosong yang infimumnya dapat diambil, dan
$\theta A$ selalu terdefinisi dalam $[0,\infty]$.
}%end of comment

Fungsi $\theta$ ini disebut {\bf ukuran luar Lebesgue} pada
$\BbbR^{r}$; frasa ini dibenarkan oleh (a) dari proposisi berikutnya.


\leader{115D}{Proposisi} (a) $\theta$ adalah ukuran luar pada
$\BbbR^r$.

(b) $\theta I=\lambda I$ untuk setiap interval setengah terbuka
$I\subseteq\BbbR^r$.

\proof{{\bf (a)(i)} $\theta$ mengambil nilai dalam $[0,\infty]$
karena setiap $\theta A$ merupakan infimum suatu subhimpunan tak kosong
dari $[0,\infty]$.

\medskip

\quad{\bf (ii)} $\theta\emptyset=0$ karena (misalnya) jika kita
menetapkan $I_j=\emptyset$ untuk setiap $j$, maka setiap $I_j$ adalah
interval setengah terbuka (menurut kesepakatan yang saya gunakan),
$\emptyset\subseteq\bigcup_{j\in\Bbb N}I_j$
dan $\sum_{j=0}^{\infty}\lambda I_j=0$.

\medskip

\quad{\bf (iii)} Jika $A\subseteq B$, maka setiap kali
$B\subseteq\bigcup_{j\in\Bbb N}I_j$ kita mempunyai
$A\subseteq\bigcup_{j\in\Bbb N}I_j$, sehingga $\theta A$ adalah
infimum suatu himpunan yang memuat himpunan yang terlibat dalam
definisi $\theta B$, dan
$\theta A\le\theta B$.

\medskip

\quad{\bf (iv)} Sekarang andaikan bahwa
$\langle A_n\rangle_{n\in\Bbb N}$ adalah barisan subhimpunan
$\BbbR^r$, dengan gabungan $A$. Untuk setiap
$\epsilon>0$, kita dapat memilih, bagi setiap $n\in\Bbb N$, suatu
barisan $\langle I_{nj}\rangle_{j\in\Bbb N}$ yang terdiri atas interval
setengah terbuka sedemikian sehingga
$A_n\subseteq\bigcup_{j\in\Bbb N}I_{nj}$ dan
$\sum_{j=0}^{\infty}\lambda I_{nj}\le\theta A_n+2^{-n}\epsilon$.
(Anda mungkin perlu memeriksa bahwa perumusan ini sah, baik ketika
$\theta A_n$ berhingga maupun tak hingga.) Sekarang, menurut 111F(b-ii),
terdapat bijeksi dari $\Bbb N$ ke $\Bbb N\times\Bbb N$; nyatakan
bijeksi ini dalam bentuk $m\mapsto(k_m,l_m)$. Maka kita memperoleh

\Centerline{$\sum_{m=0}^{\infty}\lambda I_{k_m,l_m}
=\sum_{n=0}^{\infty}\sum_{j=0}^{\infty}\lambda I_{nj}$.}

\noindent (Untuk melihat ini, perhatikan bahwa karena setiap
$\lambda I_{nj}$ lebih besar daripada atau sama dengan $0$, dan
$m\mapsto(k_m,l_m)$ adalah bijeksi, kedua jumlah sama dengan

\Centerline{$\sup_{K\subseteq\Bbb N\times\Bbb N\text{ berhingga}}
\sum_{(n,j)\in K}\lambda I_{nj}$.}

\noindent Atau lihat argumen yang dituliskan lengkap dalam 114D.) Namun
sekarang $\langle I_{k_m,l_m}\rangle_{m\in\Bbb N}$ adalah barisan
interval setengah terbuka dan

\Centerline{$A=\bigcup_{n\in\Bbb N}A_n
\subseteq\bigcup_{n\in\Bbb N}\bigcup_{j\in\Bbb N}I_{nj}
=\bigcup_{m\in\Bbb N}I_{k_m,l_m}$,}

\noindent sehingga

$$\eqalign{\theta A
&\le\sum_{m=0}^{\infty}\lambda I_{k_m,l_m}
=\sum_{n=0}^{\infty}\sum_{j=0}^{\infty}\lambda I_{nj}\cr
&\le\sum_{n=0}^{\infty}(\theta A_n+2^{-n}\epsilon)
=\sum_{n=0}^{\infty}\theta A_n+\sum_{n=0}^{\infty}2^{-n}\epsilon
=\sum_{n=0}^{\infty}\theta A_n+2\epsilon.\cr}$$

\noindent Karena $\epsilon$ sebarang, $\theta
A\le\sum_{n=0}^{\infty}\theta A_n$ (sekali lagi, Anda sebaiknya
memeriksa bahwa ini sah, baik
$\sum_{n=0}^{\infty}\theta A_n$ berhingga maupun tidak).
Karena $\langle A_n\rangle_{n\in\Bbb N}$ sebarang, $\theta$ adalah
ukuran luar.

\medskip

{\bf (b)} Karena kita selalu dapat mengambil $I_0=I$ dan
$I_j=\emptyset$ untuk $j\ge 1$, untuk memperoleh barisan interval
setengah terbuka yang menutupi $I$ dengan
$\sum_{j=0}^{\infty}\lambda I_j=\lambda I$, kita tentu mempunyai
$\theta I\le\lambda I$. Untuk ketaksamaan sebaliknya, gunakan 115B;
jika $I\subseteq\bigcup_{j\in\Bbb N}I_j$, maka $\lambda
I\le\sum_{j=0}^{\infty}\lambda I_j$; karena
$\langle I_j\rangle_{j\in\Bbb N}$ sebarang, $\theta I\ge\lambda I$ dan
$\theta I=\lambda I$, sebagaimana diperlukan.
}%end of proof of 115D


\leader{115E}{Definisi} Karena ukuran luar
Lebesgue\cmmnt{ (115C)} adalah\cmmnt{ memang}
ukuran luar\cmmnt{ (115Da)},
kita dapat menggunakannya untuk mengonstruksi suatu ukuran $\mu$, dengan
metode \Caratheodory\cmmnt{ (113C)}.
Ukuran ini adalah {\bf ukuran Lebesgue pada $\BbbR^r$}.
Himpunan-himpunan $E$ dengan $\mu E$ terdefinisi\cmmnt{ (yakni,
yang memenuhi $\theta(A\cap E)+\theta(A\setminus E)=\theta A$ untuk
setiap $A\subseteq\BbbR^r$)} disebut {\bf terukur Lebesgue}.

Himpunan-himpunan yang terabaikan untuk $\mu$ disebut {\bf terabaikan
Lebesgue}\cmmnt{; perhatikan bahwa himpunan-himpunan ini tepatlah
himpunan $A$ yang memenuhi $\theta A=0$, dan semuanya terukur Lebesgue
(113Xa)}.


\leader{115F}{Lema} Jika $i\le r$ dan $\xi\in\Bbb R$, maka
$H_{i\xi}=\{y:\eta_i<\xi\}$ terukur Lebesgue.

\proof{ Tuliskan $H$ untuk $H_{i\xi}$.

\medskip

{\bf (a)} Intinya adalah bahwa $\lambda I=\lambda(I\cap
H)+\lambda(I\setminus H)$ untuk setiap interval setengah terbuka
$I\subseteq\BbbR^r$. \Prf\ Jika $I\subseteq H$ atau
$I\cap H=\emptyset$, hal ini langsung jelas. Jika tidak, $I$ harus
berbentuk $\coint{a,b}$, dengan $\alpha_i<\xi<\beta_i$. Sekarang
$I\cap H=\coint{a,x}$ dan
$I\setminus H=\coint{y,b}$, dengan $\xi_j=\beta_j$ untuk $j\ne i$,
$\xi_i=\xi$, $\eta_j=\alpha_j$ untuk $j\ne i$, $\eta_i=\xi$, sehingga
keduanya merupakan interval setengah terbuka, dan

$$\eqalign{\lambda(I\cap H)+\lambda(I\setminus H)
&=(\xi-\alpha_i)\prod_{j\ne i}(\beta_j-\alpha_j)
+(\beta_i-\xi)\prod_{j\ne i}(\beta_j-\alpha_j)\cr
&=(\beta_i-\alpha_i)\prod_{j\ne i}(\beta_j-\alpha_j)
=\lambda I.\text{ \Qed}\cr}$$

\medskip

{\bf (b)} Sekarang andaikan bahwa $A$ adalah sembarang subhimpunan
$\BbbR^r$, dan $\epsilon>0$. Maka kita dapat menemukan barisan
$\langle I_j\rangle_{j\in\Bbb N}$
yang terdiri atas interval setengah terbuka sedemikian sehingga
$A\subseteq\bigcup_{j\in\Bbb N}I_j$ dan
$\sum_{j=0}^{\infty}\lambda I_j\le\theta A+\epsilon$. Dalam hal ini,
$\langle I_j\cap H\rangle_{j\in\Bbb N}$ dan
$\langle I_j\setminus H\rangle_{j\in\Bbb N}$ adalah barisan interval
setengah terbuka,
$A\cap H\subseteq\bigcup_{j\in\Bbb N}(I_j\cap H)$ dan
$A\setminus H\subseteq\bigcup_{j\in\Bbb N}(I_j\setminus H)$. Jadi

$$\eqalign{\theta(A\cap H)+\theta(A\setminus H)
&\le\sum_{j=0}^{\infty}\lambda(I_j\cap
H)+\sum_{j=0}^{\infty}\lambda(I_j\setminus H) \cr
&=\sum_{j=0}^{\infty}\lambda I_j
\le\theta A+\epsilon.\cr}$$

\noindent Karena $\epsilon$ sebarang,
$\theta(A\cap H)+\theta(A\setminus H)\le\theta A$;
karena $A$ sebarang, $H$ terukur, sebagaimana dicatat dalam 113D.
}%end of proof of 115F

\leader{115G}{Proposisi} Semua subhimpunan Borel dari $\BbbR^r$ terukur
Lebesgue; khususnya, semua himpunan terbuka dan semua himpunan dari
kelas-kelas berikut, beserta gabungan terhitung dari himpunan-himpunan
tersebut:

\inset{interval terbuka
$\ooint{a,b}=\{x:x\in\BbbR^r,\,
\alpha_i<\xi_i<\beta_i\Forall i\le r\}$,
dengan $-\infty\le\alpha_i<\beta_i\le\infty$ untuk setiap $i\le r$;

interval tertutup $[a,b]
=\{x:x\in\BbbR^r,\,\alpha_i\le\xi_i\le\beta_i
  \Forall i\le r\}$, dengan
$-\infty<\alpha_i<\beta_i<\infty$ untuk setiap $i\le r$.}

\noindent Kita\cmmnt{ juga} mempunyai rumus berikut untuk ukuran
himpunan-himpunan tersebut, dengan $\mu$ menyatakan ukuran Lebesgue:

\Centerline{$\mu\ooint{a,b}=\mu[a,b]=\prod_{i=1}^r\beta_i-\alpha_i$}

\noindent setiap kali $a\le b$ dalam $\BbbR^r$. Akibatnya, setiap
subhimpunan terhitung dari $\BbbR^r$ terukur dan berukuran nol.

\proof{{\bf (a)} Pertama-tama saya tunjukkan bahwa semua subhimpunan
terbuka dari $\BbbR^r$ terukur. \Prf\ Misalkan $G\subseteq\BbbR^r$
terbuka. Misalkan $K\subseteq\BbbQ^r\times\BbbQ^r$ adalah himpunan
pasangan $(c,d)$ dari tupel-$r$ bilangan rasional sedemikian sehingga
$\coint{c,d}\subseteq G$.
Sekarang, menurut catatan
dalam 111E-111F --- khususnya, 111Eb, yang menunjukkan bahwa $\Bbb Q$
terhitung, 111F(b-iii), yang menunjukkan bahwa produk dua himpunan
terhitung adalah terhitung, dan 111F(b-i), yang menunjukkan bahwa
subhimpunan dari himpunan terhitung adalah terhitung --- kita melihat,
dengan induksi pada $r$, bahwa $\BbbQ^r$ terhitung, dan bahwa $K$
terhitung. Selain itu, setiap $\coint{c,d}$ terukur, karena merupakan

\Centerline{$\bigcap_{i\le r}H_{i\delta_i}\setminus H_{i\gamma_i}$,}

\noindent dalam notasi 115F, jika $c=(\gamma_1,\ldots,\gamma_r)$ dan
$d=(\delta_1,\ldots,\delta_r)$. Jadi, menurut 111Fa,
$G'=\bigcup_{(r,s)\in K}\coint{r,s}$ terukur.

Menurut definisi $K$, $G'\subseteq G$. Di pihak lain, jika $x\in G$,
terdapat $\epsilon>0$ sedemikian sehingga $y\in G$ setiap kali
$\|y-x\|<\epsilon$.
Sekarang, untuk setiap $i$, terdapat bilangan-bilangan rasional
$\gamma_i$, $\delta_i$ sedemikian sehingga
$\gamma_i\le\xi_i<\delta_i$ dan
$\delta_i-\gamma_i\le\Bover{\epsilon}{\sqrt r}$. Jika
$y\in\coint{c,d}$, maka
$|\eta_i-\xi_i|<\Bover{\epsilon}{\sqrt r}$ untuk setiap $i$, sehingga
$\|y-x\|<\epsilon$ dan $y\in G$. Dengan demikian $(c,d)\in K$ dan
$x\in\coint{c,d}\subseteq G'$.
Karena $x$ sebarang, $G=G'$ dan $G$ terukur.\ \Qed

\medskip

{\bf (b)} Sekarang keluarga $\Sigma$ dari himpunan-himpunan terukur
Lebesgue adalah aljabar-$\sigma$ atas subhimpunan-subhimpunan $\BbbR^r$
yang memuat keluarga himpunan terbuka, sehingga harus memuat setiap
himpunan Borel, menurut definisi himpunan Borel (111G).

\medskip

{\bf (c)} Di antara jenis-jenis interval yang ditinjau, semua interval
terbuka memang merupakan himpunan terbuka, sehingga tentu merupakan
himpunan Borel. Sebuah interval tertutup $[a,b]$ dapat dinyatakan sebagai
irisan
$\bigcap_{n\in\Bbb N}\ooint{a-2^{-n}\tbf{1},b+2^{-n}\tbf{1}}$ dari suatu
barisan interval terbuka, sehingga merupakan himpunan Borel.

\medskip

{\bf (d)} Untuk menghitung ukurannya, dari 115Db kita sudah mengetahui
bahwa $\mu\coint{a,b}=\prod_{i=1}^r\beta_i-\alpha_i$ jika $a\le b$.
Untuk jenis-jenis interval berbatas lainnya, cukup diperhatikan bahwa
jika $-\infty<\alpha_i<\beta_i<\infty$ untuk setiap $i$, maka


\Centerline{$\coint{a+\epsilon\tbf{1},b}
\subseteq\ooint{a,b}\subseteq[a,b]\subseteq
\coint{a,b+\epsilon\tbf{1}}$}

\noindent setiap kali $\epsilon>0$ dalam $\Bbb R$. Jadi

\Centerline{$\mu\ooint{a,b}\le\mu[a,b]
\le\inf_{\epsilon>0}\mu\coint{a,b+\epsilon\tbf{1}}
=\inf_{\epsilon>0}\prod_{i=1}^r(\beta_i-\alpha_i+\epsilon)
=\prod_{i=1}^r\beta_i-\alpha_i$.}

\noindent Jika $\beta_i=\alpha_i$ untuk suatu $i$, maka kita harus
mempunyai

\Centerline{$\mu\ooint{a,b}=\mu[a,b]=
0=\prod_{i=1}^r\beta_i-\alpha_i$.}

\noindent Jika $\beta_i>\alpha_i$ untuk setiap $i$, tetapkan
$\epsilon_0=\min_{i\le r}\beta_i-\alpha_i>0$; maka

$$\eqalign{\mu[a,b]\ge\mu\ooint{a,b}
&\ge\sup_{0<\epsilon\le\epsilon_0}\mu\coint{a+\epsilon\tbf{1},b}\cr
&=\sup_{0<\epsilon\le\epsilon_0}
  \prod_{i=1}^r(\beta_i-\alpha_i-\epsilon)
=\prod_{i=1}^r\beta_i-\alpha_i.\cr}$$

\noindent Jadi, dalam kasus ini,

\Centerline{$\prod_{i=1}^r\beta_i-\alpha_i\le\mu\ooint{a,b}
\le\mu[a,b]\le\prod_{i=1}^r\beta_i-\alpha_i$}

\noindent dan

\Centerline{$\mu\ooint{a,b}
=\mu[a,b]=\prod_{i=1}^r\beta_i-\alpha_i$.}

\medskip

{\bf (e)} Menurut (d), $\mu\{a\}=\mu[a,a]=0$ untuk setiap $a$. Jika
$A\subseteq\BbbR^r$ terhitung, himpunan tersebut kosong atau dapat
dinyatakan sebagai $\{a_n:n\in\Bbb N\}$. Dalam kasus pertama,
$\mu A=\mu\emptyset=0$; dalam kasus kedua,
$A=\bigcup_{n\in\Bbb N}\{a_n\}$ adalah himpunan Borel dan
$\mu A\le\sum_{n=0}^{\infty}\mu\{a_n\}=0$.
}%end of proof of 115G

\exercises{
\leader{115X}{Latihan dasar} Jika Anda melewati \S114, sekarang Anda
sebaiknya kembali ke 114X dan memastikan bahwa Anda dapat mengerjakan
latihan-latihan di sana maupun latihan-latihan di bawah ini.

\header{115Xa}{\bf (a)} Tunjukkan bahwa jika $I$, $J$ adalah interval
setengah terbuka dalam $\BbbR^r$, maka $I\setminus J$ dapat dinyatakan
sebagai gabungan paling banyak $2r$
interval setengah terbuka yang saling lepas. Selanjutnya, tunjukkan bahwa
(i) setiap gabungan berhingga interval-interval setengah terbuka dapat
dinyatakan sebagai gabungan berhingga interval-interval setengah terbuka
yang saling lepas (ii) setiap gabungan terhitung interval-interval
setengah terbuka dapat dinyatakan sebagai gabungan suatu barisan
interval setengah terbuka yang saling lepas.

\sqheader 115Xb Tuliskan $\theta$ untuk ukuran luar Lebesgue, $\mu$
untuk ukuran Lebesgue pada $\BbbR^r$. Tunjukkan bahwa
$\theta A=\inf\{\mu E:E$ terukur Lebesgue, $A\subseteq E\}$ untuk setiap
$A\subseteq\BbbR^r$.
\Hint{Tinjau himpunan-himpunan $E$ berbentuk
$\bigcup_{j\in\Bbb N}I_j$, dengan $\sequence{j}{I_j}$ suatu barisan
interval setengah terbuka.}

\header{115Xc}{\bf (c)} Misalkan $E\subseteq\BbbR^r$ suatu himpunan
yang mempunyai ukuran berhingga menurut ukuran Lebesgue $\mu$.
Tunjukkan bahwa untuk setiap $\epsilon>0$ terdapat keluarga saling lepas
$I_0,\ldots,I_n$ yang terdiri atas interval setengah terbuka sedemikian
sehingga $\mu(E\symmdiff\bigcup_{j\le n}I_j)\le\epsilon$. \Hint{misalkan
$\langle J_j\rangle_{j\in\Bbb N}$ suatu barisan interval setengah
terbuka sedemikian sehingga
$E\subseteq\bigcup_{j\in\Bbb N}J_j$ dan
$\sum_{j=0}^{\infty}\mu J_j\le\mu E+\bover{1}{2}\epsilon$. Sekarang
ambil $m$ yang cukup besar dan nyatakan $\bigcup_{j\le m}J_j$ sebagai
gabungan saling lepas interval-interval setengah terbuka.}

\sqheader 115Xd Andaikan bahwa $c\in\BbbR^r$. (i) Tunjukkan bahwa
$\theta(A+c)=\theta A$ untuk setiap $A\subseteq\BbbR^r$, dengan
$A+c=\{x+c:x\in A\}$. (ii)
Tunjukkan bahwa jika $E\subseteq \BbbR^r$ terukur, maka $E+c$ juga
terukur, dan dalam hal ini $\mu(E+c)=\mu E$.

\header{115Xe}{\bf (e)} Andaikan bahwa $\gamma>0$. (i) Tunjukkan bahwa
$\theta(\gamma A)=\gamma^r\theta A$ untuk setiap $A\subseteq\BbbR^r$,
dengan $\gamma A=\{\gamma x:x\in A\}$. (ii) Tunjukkan bahwa jika
$E\subseteq \BbbR^r$ terukur, maka $\gamma E$ juga terukur, dan dalam
hal ini $\mu(\gamma E)=\gamma^r\mu E$

\vleader{48pt}{115Y}{Latihan lanjutan (a)}
%\spheader 115Ya
(i) Andaikan bahwa $M$ adalah bilangan bulat positif dan $k_i$, $l_i$
adalah bilangan bulat untuk $1\le i\le r$. Tetapkan
$\alpha_i=k_i/M$ dan
$\beta_i=l_i/M$ untuk setiap $i$, dan $I=\coint{a,b}$. Tunjukkan bahwa
$\lambda I=\#(J)/M^r$, dengan $J$ menyatakan $\{z:z\in\BbbZ^r$,
$\bover1Mz\in I\}$. (ii) Tunjukkan bahwa jika suatu interval setengah
terbuka $I\subseteq\BbbR^r$ ditutupi oleh suatu barisan {\it berhingga}
$I_0,\ldots,I_m$ yang terdiri atas interval setengah terbuka, dan semua
koordinat yang terlibat dalam menentukan interval-interval $I$,
$I_0,\ldots,I_m$ adalah rasional, maka
$\lambda I\le\sum_{j=0}^m\lambda I_j$. (iii) Dengan mengandaikan
teorema Heine--Borel dalam bentuk

\inset{\noindent
setiap kali $[a,b]$ merupakan interval tertutup dalam $\BbbR^r$ yang
ditutupi oleh suatu barisan $\sequence{j}{\ooint{a^{(j)},b^{(j)}}}$
yang terdiri atas interval terbuka, terdapat $m\in\Bbb N$ sedemikian
sehingga
$[a,b]\subseteq\bigcup_{j\le m}\ooint{a^{(j)},b^{(j)}}$,}

\noindent buktikan 115B. \Hint{jika
$\coint{a,b}\subseteq\bigcup_{j\in\Bbb N}\coint{a^{(j)},b^{(j)}}$,
gantikan $\coint{a,b}$ dengan interval tertutup yang lebih kecil dan
setiap $\coint{a^{(j)},b^{(j)}}$ dengan interval terbuka yang lebih besar,
dengan mengubah masing-masing volume sebesar nilai yang cukup kecil.}
%{\smc Facenda \& Freniche 02}

\spheader 115Yb(i)
Tunjukkan bahwa jika $A\subseteq\BbbR^r$ dan $\epsilon>0$, terdapat
himpunan terbuka $G\supseteq A$ sedemikian sehingga
$\theta G\le\theta A+\epsilon$, dengan $\theta$ ukuran luar Lebesgue.
(ii) Tunjukkan bahwa jika $E\subseteq\BbbR^r$ terukur Lebesgue dan
$\epsilon>0$, terdapat himpunan terbuka $G\supseteq E$ sedemikian
sehingga $\mu(G\setminus E)\le\epsilon$, dengan $\mu$ ukuran Lebesgue.
\Hint{tinjau terlebih dahulu kasus $E$ berbatas.} (iii) Tunjukkan bahwa
jika $E\subseteq\BbbR^r$ terukur Lebesgue, terdapat himpunan-himpunan
Borel $H_1$, $H_2$ sedemikian sehingga $H_1\subseteq E\subseteq H_2$
dan $\mu(H_2\setminus E)=\mu(E\setminus H_1)=0$.
\Hint{gunakan (ii) untuk menemukan $H_2$, lalu tinjau komplemen
$E$.}

\spheader 115Yc Tuliskan $\theta$ untuk ukuran luar Lebesgue pada
$\BbbR^r$. Tunjukkan
bahwa suatu himpunan $E\subseteq\BbbR^r$ terukur Lebesgue jika dan hanya
jika
$\theta([-\tbf{n},\tbf{n}]\cap E)+\theta([-\tbf{n},\tbf{n}]\setminus
E)=(2n)^r$ untuk setiap $n\in\Bbb N$, dengan
$\tbf{n}=(n,\ldots,n)$. \Hint{gunakan 115Yb untuk menunjukkan bahwa
bagi setiap $n$ terdapat
himpunan-himpunan terukur $F_n$, $H_n$ sedemikian sehingga
$F_n\subseteq[-\tbf{n},\tbf{n}]\cap
E\subseteq H_n$ dan $H_n\setminus F_n$ terabaikan.}

\header{115Yd}{\bf (d)} Dengan mengandaikan bahwa terdapat himpunan
$A\subseteq\Bbb R$ yang bukan himpunan Borel, tunjukkan bahwa terdapat
keluarga $\Cal E$ yang terdiri atas interval setengah terbuka dalam
$\BbbR^2$ sedemikian sehingga $\bigcup\Cal E$ bukan himpunan Borel.
\Hint{tinjau $\Cal
E=\{\coint{\xi,1+\xi}\times\coint{-\xi,1-\xi}:\xi\in A\}$.}

\header{115Ye}{\bf (e)} Misalkan $X$ suatu himpunan dan $\Cal A$ suatu
{\bf semiring} yang terdiri atas subhimpunan-subhimpunan $X$, yakni,
suatu keluarga subhimpunan $X$ sedemikian sehingga

\inset{$\emptyset\in\Cal A$,}

\inset{$E\cap F\in\Cal A$ untuk semua $E$, $F\in\Cal A$,}

\inset{setiap kali $E$, $F\in\Cal A$, terdapat himpunan-himpunan saling
lepas $E_0,\ldots,E_n\in\Cal A$ sedemikian sehingga
$E\setminus F=E_0\cup\ldots\cup E_n$.}

\noindent Misalkan $\lambda:\Cal A\to[0,\infty]$ suatu fungsional
sedemikian sehingga

\inset{$\lambda\emptyset=0$,}

\inset{$\lambda E=\sum_{i=0}^{\infty}\lambda E_i$ setiap kali
$E\in\Cal A$ dan $\sequence{i}{E_i}$ adalah barisan saling lepas dalam
$\Cal A$ dengan gabungan $E$.}

\noindent Tunjukkan bahwa terdapat suatu ukuran $\mu$ pada $X$ yang
memperluas $\lambda$. \Hint{gunakan metode 113Yi.}
}%end of exercises

\endnotes{
\Notesheader{115} Dalam catatan untuk \S114, saya meninjau sekilas
metode-metode yang sejauh ini tersedia bagi kita untuk mengonstruksi
ruang ukur. Sekarang ukuran Lebesgue pada $\BbbR^r$ dapat kita tambahkan
ke dalam daftar tersebut.

Jika Anda melihat kembali \S114, Anda akan melihat bahwa saya sengaja
menyalin pemaparan di sana. Saya berharap pengulangan ini membantu Anda
melihat unsur-unsur hakiki dari metode tersebut, yang berjumlah tiga:
suatu konsep primitif tentang volume (114A/115A); subaditivitas
terhitung (114B/115B); dan keterukuran blok-blok pembangun (114F/115F).

Mengenai `konsep primitif tentang volume', tidak banyak yang perlu
dikatakan. Gagasan panjang suatu interval, luas suatu persegi panjang,
dan volume suatu balok dapat ditelusuri sampai ke awal sejarah
matematika. Saya menggunakan `interval setengah terbuka', sebagaimana
didefinisikan dalam 114Aa/115Ab, semata-mata karena alasan teknis, yakni
karena interval-interval tersebut dapat dirangkai dengan rapi (lihat
115Xa dan 115Ye); jika kita memulai dengan interval `terbuka' atau
`tertutup', metode ini tetap akan berhasil. Satu hal barangkali layak
disebutkan: semua balok yang saya gunakan berorientasi tegak, dengan
rusuk-rusuk sejajar terhadap sumbu-sumbu koordinat. Sesungguhnya,
membuktikan bahwa
suatu balok dalam orientasi lain mempunyai ukuran Lebesgue yang tepat
merupakan latihan yang taktrivial, dan saya menundanya hingga Bab 26.
Untuk saat ini, kita mencari jalan aman terpendek menuju definisi yang
tepat, dan fakta bahwa merotasi suatu himpunan tidak mengubah ukuran
Lebesgue-nya harus menunggu.

Langkah besarnya ialah `subaditivitas terhitung': fakta bahwa jika suatu
balok ditutupi oleh barisan balok lain, volumenya lebih kecil daripada
atau sama dengan jumlah volume balok-balok tersebut. Hal ini tentunya
diperlukan jika balok-balok itu harus terukur dan mempunyai ukuran yang tepat,
menurut 112Cd. (Yang luar biasa ialah bahwa syarat ini nyaris sudah
mencukupi.) Di sini kita mempunyai pekerjaan yang perlu dilakukan, dan
dalam kasus berdimensi $r$ terdapat bukit terjal yang harus didaki. Anda
dapat mendakinya dalam dua tahap jika Anda mencari teorema Heine--Borel
(115Ya); tetapi sebagaimana saya coba jelaskan dalam catatan setelah
115B, menurut saya jalur ini tidak menghindari satu pun kesulitan yang
sesungguhnya.

Hal ketiga yang harus kita periksa ialah bahwa balok-balok tersebut
terukur dalam pengertian teknis yang diuraikan oleh teorema
\Caratheodory. Ini karena balok-balok itu dapat diperoleh dari
setengah-ruang melalui operasi irisan, gabungan, dan pengambilan
komplemen, sedangkan setengah-ruang terukur karena alasan yang sangat
langsung (114F/115F). Setelah itu kita sudah melangkah jauh, dan saya
hanya melakukan sedikit lagi: sekadar memeriksa bahwa himpunan terbuka,
dan dengan demikian himpunan Borel, terukur, serta bahwa interval
tertutup dan terbuka mempunyai ukuran yang tepat (114G/115G). Beberapa
sifat ukuran Lebesgue yang lain dapat ditemukan dalam \S134. Namun,
setiap jilid---bahkan hampir setiap bab---dalam risalah ini akan
memperkenalkan ciri-ciri lebih lanjut dari konstruksi yang luar biasa
ini.
}%end of notes

\frnewpage
